\chapter*{Introduction}
\addcontentsline{toc}{chapter}{Introduction}

Le stage de perfectionnement représente une étape cruciale dans la formation de l'étudiant, lui permettant de confronter ses acquis théoriques aux réalités du monde professionnel. C'est dans ce cadre que j'ai eu l'opportunité d'effectuer mon stage au sein de la société SCPC Comptoir Hammami.

L'objectif de ce stage était de concevoir et développer une application web, intitulée "Central de Réclamation", destinée à moderniser et optimiser la gestion des réclamations clients, qui était jusqu'alors traitée manuellement.

Ce rapport retrace les différentes étapes de ce projet, de l'étude de l'existant à la réalisation finale. Il est structuré comme suit :
\begin{itemize}
    \item \textbf{Chapitre 1 : Etude de l'organisme.} Présentation de la société SCPC Comptoir Hammami et de son environnement.
    \item \textbf{Chapitre 2 : Présentation du thème.} Définition du périmètre du projet et planning prévisionnel.
    \item \textbf{Chapitre 3 : Infrastructure informatique.} Description des ressources matérielles et logicielles.
    \item \textbf{Chapitre 4 : Etude préalable.} Analyse de l'existant, critique et justification de la solution.
    \item \textbf{Chapitre 5 : Etude conceptuelle.} Conception UML (Cas d'utilisation, Séquence, Classes).
    \item \textbf{Chapitre 6 : Réalisation.} Présentation de l'architecture, de la sécurité et des interfaces finales.
\end{itemize}
